
\section{模型检验}

本节对前文建立的主要模型进行检验,包括误差分析、交叉验证和灵敏度分析三个方面。

\subsection{问题一模型的检验}

对于问题一的【主方法】,采用【5 折 / 10 折】交叉验证评估其泛化能力。结果显示,【指标】为【值 ± 标准差】,【进一步说明】。

\subsection{问题二模型的检验}

对问题二的【主方法】,采用交叉验证和【误差分析方法】两种方式进行检验。

【描述检验方法与结果】。

\subsection{问题三模型的检验}

对问题三的【主方法】,从以下【N】个角度进行检验。

\textbf{交叉验证:}【结果说明】。

\textbf{【其他检验方法】:}【结果说明】。

\textbf{【灵敏度分析】:}对【关键因素】进行 $\pm 20\%$ 缩放,观察【响应指标】变化。结果显示,【最敏感因素】对【结果】影响最显著。

\subsection{问题四模型的检验}

对问题四的【分类/预测】模型,主要采用 5 折分层交叉验证和【其他评估方法】进行检验。

\textbf{5 折交叉验证:}【主方法】的交叉验证准确率【值】, F1 分数【值】;【对比方法】的交叉验证准确率【值】, F1 分数【值】。

\textbf{【其他方法】:}【结果说明】。

\textbf{类别不平衡处理:}【说明处理方法及效果】。

综合上述检验,本文所建立的主要模型在交叉验证和误差传播分析中均表现出较好的稳健性,【进一步说明】。
